\begin{tikzpicture}[font=\sffamily]
  % Izquierda: Mapa de Valor (Cuadrado)
  \fill[cepanelcrimson] (0, 0) rectangle (5.6, 5.6);
  \draw[cecrimson, line width=1.1pt] (0, 0) rectangle (5.6, 5.6);
  \draw[cecrimson!70, line width=0.8pt] (0, 0) -- (2.8, 2.8) -- (0, 5.6);
  \draw[cecrimson!70, line width=0.8pt] (2.8, 2.8) -- (5.6, 2.8);

  \node[font=\sffamily\bfseries\small, cecrimson] at (2.8, 6.05) {MAPA DE VALOR (Oferta)};
  \node[text width=1.95cm, align=center, font=\sffamily\scriptsize] at (1.02, 2.8)
    {\textcolor{cecrimson}{\bfseries Productos y\\servicios}\\[3pt]\color{ceslate} Lista de componentes físicos, digitales o financieros.};
  \node[text width=2.45cm, align=center, font=\sffamily\scriptsize] at (3.9, 4.3)
    {\textcolor{cecrimson}{\bfseries Creadores de\\alegrías}\\[3pt]\color{ceslate} Cómo produce beneficios y resultados deseados.};
  \node[text width=2.45cm, align=center, font=\sffamily\scriptsize] at (3.9, 1.3)
    {\textcolor{cecrimson}{\bfseries Aliviadores de\\frustraciones}\\[3pt]\color{ceslate} Cómo elimina costes, riesgos o fricciones.};

  % Centro: Flechas de Encaje (FIT)
  \draw[<->, >=Stealth, line width=1.6pt, ceindigo] (5.8, 2.8) -- (7.55, 2.8);
  \node[rectangle, rounded corners=3pt, fill=ceindigo, text=white, font=\sffamily\bfseries\footnotesize, inner sep=4pt]
    at (6.675, 3.45) {ENCAJE};

  % Derecha: Perfil del Cliente (Círculo)
  \begin{scope}[shift={(10.6, 2.8)}]
    \fill[cepanelteal] (0, 0) circle (2.8cm);
    \draw[ceteal, line width=1.1pt] (0, 0) circle (2.8cm);
    \draw[ceteal!70, line width=0.8pt] (-2.8, 0) -- (0, 0);
    \draw[ceteal!70, line width=0.8pt] (0, 0) -- (1.98, 1.98);
    \draw[ceteal!70, line width=0.8pt] (0, 0) -- (1.98, -1.98);

    \node[font=\sffamily\bfseries\small, ceteal!90!black] at (0, 3.25) {PERFIL DEL CLIENTE (Demanda)};
    \node[text width=2.2cm, align=center, font=\sffamily\scriptsize] at (1.55, 0)
      {\textcolor{ceteal!90!black}{\bfseries Trabajos del\\cliente}\\[2pt]\color{ceslate} Tareas funcionales, sociales y emocionales.};
    \node[text width=2.1cm, align=center, font=\sffamily\scriptsize] at (-0.85, 1.35)
      {\textcolor{ceteal!90!black}{\bfseries Alegrías}\\[2pt]\color{ceslate} Resultados y beneficios esperados.};
    \node[text width=2.1cm, align=center, font=\sffamily\scriptsize] at (-0.85, -1.35)
      {\textcolor{ceteal!90!black}{\bfseries Frustraciones}\\[2pt]\color{ceslate} Riesgos, trabas y costes negativos.};
  \end{scope}

  % Pie explicativo de los 3 niveles de encaje
  \node[rectangle, rounded corners=3pt, draw=ceindigo!60, fill=cecream, text width=13.4cm, align=center, inner sep=5pt, font=\sffamily\scriptsize]
    at (6.7, -0.75)
    {\textbf{Tres niveles de encaje (Osterwalder et al., 2014):} (1) \textit{Sobre papel} (hipótesis Problema--Solución) $\rightarrow$ (2) \textit{En mercado} (Producto--Mercado: pagan y repiten) $\rightarrow$ (3) \textit{En escala} (Modelo rentable).};
\end{tikzpicture}